% GENERATED FILE. Edit canonical lessons, then run npm run generate:books.
% canonical-source-hash: fnv1a64:36d601ea994237a3
% canonical-lessons: MR-C538-tal, MR-C538-nartak, MR-C538-ghungru, MR-C538-theka-dharne, MR-C538-dolne

\chapter{Rhythm and Dance}
\label{ch:mr-rhythm-and-dance}
\hlchaptermodality{538}

\begin{chapteropening}
\textbf{By the end of this chapter:} \emph{I can talk about rhythm, dancers and moving to music.}

Complete the last lesson of chapter 538: I can talk about rhythm, dancers and moving to music.
\end{chapteropening}

\section[\texorpdfstring{\mr{ताल} (tāl)}{ताल (tāl)}]{\mr{ताल} (tāl) — rhythm; a beat}

\label{lesson:MR-C538-tal}



\begin{quote}
Before the new one: say the Marathi for a costume, then the Marathi for to portray.
\end{quote}

\subsection*{You'll want to know: \mr{ताल}}
\textbf{\mr{ताल}} — \emph{tāl} — \textquotedblleft{}rhythm; a beat\textquotedblright{}.

\begin{quote}
\mr{ढोलाचा} \mr{ताल} \mr{ऐकून} \mr{सगळे} \mr{नाचू} \mr{लागले}.
\end{quote}

\begin{quote}
\emph{ḍholācā tāl aikūn sagḷe nācū lāgle.}
\end{quote}

\begin{quote}
Hearing the drum's rhythm, everyone began to dance.
\end{quote}

Masculine.

\begin{cousinweb}
\textbf{\mr{ताल}} (\emph{tāl}) is a tatsama, taken from Sanskrit \emph{tāla}, \textquotedblleft{}clapping the hands, musical time\textquotedblright{} (Monier-Williams).
\end{cousinweb}

\begin{grammarlens}[title={What you've built}]
One more word for describing what you lived and felt.
\end{grammarlens}

\subsection*{Your turn}
\begin{itemize}
\raggedright
  \item \emph{Say it:} \emph{tāl}
  \item \emph{Say it:} \emph{tāl}, once more
  \item \emph{Say it:} \emph{ḍholācā tāl aikūn sagḷe nācū lāgle.} — Hearing the drum's rhythm, everyone began to dance.
  \item \emph{Recall:} say the Marathi for a costume, then the Marathi for to portray, then say \emph{tāl} again
\end{itemize}

\begin{tcolorbox}[breakable,colback=teal!4,colframe=teal!35!black,title={Before you move on}]
What does \mr{ताल} mean? (\textquotedblleft{}Rhythm; a beat\textquotedblright{}.) Say it once more.
\end{tcolorbox}
\section[\texorpdfstring{\mr{नर्तक} (nartak)}{नर्तक (nartak)}]{\mr{नर्तक} (nartak) — a dancer}

\label{lesson:MR-C538-nartak}



\begin{quote}
Before the new one: say the Marathi for to portray, then the Marathi for rhythm.
\end{quote}

\subsection*{You'll want to know: \mr{नर्तक}}
\textbf{\mr{नर्तक}} — \emph{nartak} — \textquotedblleft{}a dancer\textquotedblright{}.

\begin{quote}
\mr{रंगमंचावर} \mr{दहा} \mr{नर्तक} \mr{नाचत} \mr{होते}.
\end{quote}

\begin{quote}
\emph{raṅgmañcāvar dahā nartak nācat hote.}
\end{quote}

\begin{quote}
Ten dancers were dancing on the stage.
\end{quote}

Masculine.

A woman dancer is \mr{नर्तकी}.

\begin{cousinweb}
\textbf{\mr{नर्तक}} (\emph{nartak}) is a tatsama, taken from Sanskrit \emph{nartaka}, \textquotedblleft{}a dancer\textquotedblright{} (Monier-Williams), from the root of \mr{नृत्य}.
\end{cousinweb}

\begin{grammarlens}[title={What you've built}]
One more word for describing what you lived and felt.
\end{grammarlens}

\subsection*{Your turn}
\begin{itemize}
\raggedright
  \item \emph{Say it:} \emph{nartak}
  \item \emph{Say it:} \emph{nartak}, once more
  \item \emph{Say it:} \emph{raṅgmañcāvar dahā nartak nācat hote.} — Ten dancers were dancing on the stage.
  \item \emph{Recall:} say the Marathi for to portray, then the Marathi for rhythm, then say \emph{nartak} again
\end{itemize}

\begin{tcolorbox}[breakable,colback=teal!4,colframe=teal!35!black,title={Before you move on}]
What does \mr{नर्तक} mean? (\textquotedblleft{}A dancer\textquotedblright{}.) Say it once more.
\end{tcolorbox}
\section[\texorpdfstring{\mr{घुंगरू} (ghuṅgrū)}{घुंगरू (ghuṅgrū)}]{\mr{घुंगरू} (ghuṅgrū) — ankle bells}

\label{lesson:MR-C538-ghungru}



\begin{quote}
Before the new one: say the Marathi for rhythm, then the Marathi for a dancer.
\end{quote}

\subsection*{You'll want to know: \mr{घुंगरू}}
\textbf{\mr{घुंगरू}} — \emph{ghuṅgrū} — \textquotedblleft{}ankle bells\textquotedblright{}.

\begin{quote}
\mr{नाचताना} \mr{तिच्या} \mr{घुंगरांचा} \mr{आवाज} \mr{येत} \mr{होता}.
\end{quote}

\begin{quote}
\emph{nāctānā ticyā ghuṅgrāñcā āvāj yet hotā.}
\end{quote}

\begin{quote}
As she danced the sound of her ankle bells could be heard.
\end{quote}

Neuter.

\begin{cousinweb}
\textbf{\mr{घुंगरू}} (\emph{ghuṅgrū}) is inherited: CDIAL lists Marathi \emph{ghũghrū}, \textquotedblleft{}bell\textquotedblright{}, under \emph{*ghuṅghura} (CDIAL 4477), a word that copies the ringing.
\end{cousinweb}

\begin{grammarlens}[title={What you've built}]
One more word for describing what you lived and felt.
\end{grammarlens}

\subsection*{Your turn}
\begin{itemize}
\raggedright
  \item \emph{Say it:} \emph{ghuṅgrū}
  \item \emph{Say it:} \emph{ghuṅgrū}, once more
  \item \emph{Say it:} \emph{nāctānā ticyā ghuṅgrāñcā āvāj yet hotā.} — As she danced the sound of her ankle bells could be heard.
  \item \emph{Recall:} say the Marathi for rhythm, then the Marathi for a dancer, then say \emph{ghuṅgrū} again
\end{itemize}

\begin{tcolorbox}[breakable,colback=teal!4,colframe=teal!35!black,title={Before you move on}]
What does \mr{घुंगरू} mean? (\textquotedblleft{}Ankle bells\textquotedblright{}.) Say it once more.
\end{tcolorbox}
\section[\texorpdfstring{\mr{ठेका} \mr{धरणे} (ṭhekā dharṇe)}{ठेका धरणे (ṭhekā dharṇe)}]{\mr{ठेका} \mr{धरणे} (ṭhekā dharṇe) — to keep the beat}

\label{lesson:MR-C538-theka-dharne}



\begin{quote}
Before the new one: say the Marathi for a dancer, then the Marathi for ankle bells.
\end{quote}

\subsection*{You'll want to know: \mr{ठेका} \mr{धरणे}}
\textbf{\mr{ठेका} \mr{धरणे}} — \emph{ṭhekā dharṇe} — \textquotedblleft{}to keep the beat\textquotedblright{}.

\begin{quote}
\mr{गाणं} \mr{सुरू} \mr{झालं} \mr{आणि} \mr{सगळ्यांनी} \mr{टाळ्यांनी} \mr{ठेका} \mr{धरला}.
\end{quote}

\begin{quote}
\emph{gāṇaṁ surū jhālaṁ āṇi sagḷyānnī ṭāḷyānnī ṭhekā dharlā.}
\end{quote}

\begin{quote}
The song began and everyone kept the beat with their claps.
\end{quote}

Past, with -\mr{ने} (he …): \mr{त्याने} \mr{ठेका} \mr{धरला}, \emph{tyāne ṭhekā dharlā}.

\begin{cousinweb}
\textbf{\mr{ठेका} \mr{धरणे}} (\emph{ṭhekā dharṇe}) is \mr{ठेका}, inherited: CDIAL lists Marathi \emph{ṭhekā}, \textquotedblleft{}mode of beating a drum\textquotedblright{}, under \emph{*ṭhēkk}, \textquotedblleft{}fix\textquotedblright{} (CDIAL 5509), with \mr{धरणे}, \textquotedblleft{}to hold\textquotedblright{}.
\end{cousinweb}

\begin{grammarlens}[title={What you've built}]
One more word for describing what you lived and felt.
\end{grammarlens}

\subsection*{Your turn}
\begin{itemize}
\raggedright
  \item \emph{Say it:} \emph{ṭhekā dharṇe}
  \item \emph{Say it:} \emph{ṭhekā dharṇe}, once more
  \item \emph{Say it:} \emph{gāṇaṁ surū jhālaṁ āṇi sagḷyānnī ṭāḷyānnī ṭhekā dharlā.} — The song began and everyone kept the beat with their claps.
  \item \emph{Recall:} say the Marathi for a dancer, then the Marathi for ankle bells, then say \emph{ṭhekā dharṇe} again
\end{itemize}

\begin{tcolorbox}[breakable,colback=teal!4,colframe=teal!35!black,title={Before you move on}]
What does \mr{ठेका} \mr{धरणे} mean? (\textquotedblleft{}To keep the beat\textquotedblright{}.) Say it once more.
\end{tcolorbox}
\section[\texorpdfstring{\mr{डोलणे} (ḍolṇe)}{डोलणे (ḍolṇe)}]{\mr{डोलणे} (ḍolṇe) — to sway}

\label{lesson:MR-C538-dolne}



\begin{quote}
Before the new one: say the Marathi for ankle bells, then the Marathi for to keep the beat.
\end{quote}

\subsection*{You'll want to know: \mr{डोलणे}}
\textbf{\mr{डोलणे}} — \emph{ḍolṇe} — \textquotedblleft{}to sway\textquotedblright{}.

\begin{quote}
\mr{संगीत} \mr{ऐकून} \mr{आजोबा} \mr{खुर्चीत} \mr{डोलू} \mr{लागले}.
\end{quote}

\begin{quote}
\emph{saṅgīt aikūn ājobā khurcīt ḍolū lāgle.}
\end{quote}

\begin{quote}
Listening to the music, Grandfather began to sway in his chair.
\end{quote}

Past (he …): \mr{तो} \mr{डोलला}, \emph{to ḍollā}.

\begin{cousinweb}
\textbf{\mr{डोलणे}} (\emph{ḍolṇe}) is perhaps kin to Sanskrit \emph{dōla}, \textquotedblleft{}swinging\textquotedblright{}, under which CDIAL lists Hindi \emph{ḍol}, \textquotedblleft{}swinging\textquotedblright{} (CDIAL 6582); CDIAL gives no Marathi verb there, so the link is not settled.
\end{cousinweb}

\begin{grammarlens}[title={What you've built}]
One more word for describing what you lived and felt.
\end{grammarlens}

\subsection*{Your turn}
\begin{itemize}
\raggedright
  \item \emph{Say it:} \emph{ḍolṇe}
  \item \emph{Say it:} \emph{ḍolṇe}, once more
  \item \emph{Say it:} \emph{saṅgīt aikūn ājobā khurcīt ḍolū lāgle.} — Listening to the music, Grandfather began to sway in his chair.
  \item \emph{Recall:} say the Marathi for ankle bells, then the Marathi for to keep the beat, then say \emph{ḍolṇe} again
\end{itemize}

\begin{tcolorbox}[breakable,colback=teal!4,colframe=teal!35!black,title={Before you move on}]
What does \mr{डोलणे} mean? (\textquotedblleft{}To sway\textquotedblright{}.) Say it once more.
\end{tcolorbox}
